\item Una fibra óptica tiene un índice de refracción $n$ y un diámetro $d$, y está rodeada por el vacío. Se envía luz por la fibra a lo largo de su eje, como se ve en la figura P34.31.
\begin{enumerate}
    \item Encuentre el mínimo radio exterior $R_{\text{mín}}$ permitido para una curva en la fibra si no ha de escapar luz.
    \item \textbf{¿Qué pasaría si?} ¿El resultado del inciso (1) pronostica un comportamiento razonable cuando $d$ se aproxima a cero? Explique.
    \item ¿Qué sucede con el radio mínimo cuando $n$ aumenta?
    \item ¿Qué sucede cuando $n$ se aproxima a 1?
    \item Evalúe $R_{\text{mín}}$ si el diámetro de la fibra es $100\ \mu\text{m}$ y su índice de refracción es $1.40$.
\end{enumerate}